\chapter{परिमाप और क्षेत्रफल}\label{ch:g5:areas}

पिछले साल हमने जाल के खाने गिने थे (\cref{ch:g4:measure}); यह
अध्याय \omterm{def:g4:measure:area}{क्षेत्रफल} का पहला \emph{सूत्र} कमाता है --- आयत के लिए
लंबाई गुना चौड़ाई --- और cm$^2$ तथा m$^2$ मात्रक सिखाता है।
त्रिभुजों और दूसरी आकृतियों के सूत्र \cref{ch:g6:measure} और
\cref{ch:g7:areas} में पकते हैं।

\section{दो अलग-अलग माप}

\begin{example}[सीमा बनाम सतह]\label{ex:g5:areas:contrast}
माली को बाड़ ख़रीदने के लिए \emph{\omterm{def:g4:measure:perimeter}{परिमाप}} चाहिए और घास का बीज
ख़रीदने के लिए \emph{\omterm{def:g4:measure:area}{क्षेत्रफल}}। दोनों एक-दूसरे के पीछे नहीं
चलते: किसी आयत को पतला और लंबा खींचने पर उसका \omterm{def:g4:measure:area}{क्षेत्रफल} वही रह
सकता है जबकि \omterm{def:g4:measure:perimeter}{परिमाप} बढ़ जाता है --- $6 \times 2$ और $4 \times 3$
के आयत मिलाकर देखो (\omterm{def:g4:measure:area}{क्षेत्रफल} दोनों का $12$, \omterm{def:g4:measure:perimeter}{परिमाप} $16$ और $14$)।
\end{example}

\section{आयत का सूत्र}

\begin{proposition}[आयत का क्षेत्रफल]\label{prop:g5:areas:rectangle}
लंबाई $L$ और चौड़ाई $w$ (एक ही मात्रक में) वाले आयत का \omterm{def:g4:measure:area}{क्षेत्रफल}
\[
A = L \times w ,
\]
वर्ग मात्रकों में होता है: cm$^2$ ($1$ cm भुजा वाले वर्ग),
m$^2$ ($1$ m भुजा वाले वर्ग), \dots
\end{proposition}

\begin{proof}[गिनकर, क्यों]
आयत को इकाई वर्गों से ढको: $L$-$L$ वर्गों की $w$ पंक्तियाँ, यानी
कुल $L \times w$ वर्ग --- ठीक वही गुणा वाला आयत जो
\cref{ch:g3:mult} में था।
\end{proof}

\begin{omfigure}
\begin{tikzpicture}[scale=0.62]
  \draw[gray!50, thin] (0,0) grid (6,3);
  \draw[very thick, omDef] (0,0) rectangle (6,3);
  \fill[omThm!30] (0,2) rectangle (1,3);
  \draw[omThm, thick] (0,2) rectangle (1,3);
  \node[right, font=\small, omThm] at (1.15,2.5) {$1$ cm$^2$};
  \node[below, font=\small] at (3,-0.15) {$6$ cm};
  \node[left, font=\small] at (0,1.5) {$3$ cm};
  \node[font=\small] at (3.9,1.2) {$6 \times 3 = 18$ cm$^2$};
\end{tikzpicture}

{\small छह-छह सेंटीमीटर-वर्गों की तीन पंक्तियाँ: सूत्र
$L \times w$ वही गिनती है, बस हमेशा के लिए लिख दी गई।}
\end{omfigure}

\begin{example}\label{ex:g5:areas:rectangle}
एक क़ालीन $2.5$ m गुणा $2$ m का है: \omterm{def:g4:measure:area}{क्षेत्रफल}
$2.5 \times 2 = 5$ m$^2$। एक डाक-टिकट $3$ cm गुणा $2.4$ cm का
है: \omterm{def:g4:measure:area}{क्षेत्रफल} $3 \times 2.4 = 7.2$ cm$^2$। सूत्र वही, मात्रक कोई
भी --- बस दोनों भुजाओं का मात्रक \emph{एक ही} हो।
\end{example}

\begin{example}[वर्ग मात्रक सौ के हिसाब से बदलते हैं]\label{ex:g5:areas:units}
$1$ m $= 100$ cm, पर $1$ m$^2 = 100 \times 100 = 10\,000$ cm$^2$:
एक वर्ग मीटर, वर्ग सेंटीमीटरों का $100$ गुणा $100$ का जाल है।
\omterm{def:g4:measure:area}{क्षेत्रफल} के मात्रक सौ-सौ की छलाँग लगाते हैं, दस-दस की नहीं।
\end{example}

\section{जोड़-तोड़ वाली आकृतियाँ}

\begin{method}[काटो, जोड़ो, घटाओ]\label{met:g5:areas:composite}
आयतों से बनी किसी आकृति (L, T, एक चौखटा) के लिए:
\begin{enumerate}
  \item उसे आयतों में काटो, या उसे पूरा करके एक बड़ा आयत बना लो;
  \item हर आयताकार \omterm{def:g4:measure:area}{क्षेत्रफल} सूत्र से निकालो;
  \item टुकड़े जोड़ दो --- या बड़े आयत में से छेद घटा दो;
  \item जाल के खानों की मोटी-मोटी गिनती से जाँच लो।
\end{enumerate}
\end{method}

\begin{example}\label{ex:g5:areas:composite}
L-आकार का एक कमरा: $6$ m $\times$ $4$ m का आयत, जिसका
$2$ m $\times$ $2$ m का कोना ग़ायब है।
\[
\text{क्षेत्रफल} = 6 \times 4 - 2 \times 2 = 24 - 4 = 20 \text{ m}^2 .
\]
या L को $6 \times 2$ की और $4 \times 2$ की पट्टियों में काट लो:
$12 + 8 = 20$ m$^2$ --- दो रास्ते, एक ही उत्तर।
\end{example}

\begin{example}[आयत का आधा]\label{ex:g5:areas:halfrect}
किसी आयत को विकर्ण पर काटने से बराबर \omterm{def:g4:measure:area}{क्षेत्रफल} वाले दो त्रिभुज
बनते हैं: हर एक आयत का \emph{\omterm{def:g3:division:half}{आधा}}। इसलिए $6$ cm और $4$ cm
भुजाओं वाले \omterm{def:g3:shapes:rightangle}{समकोण} त्रिभुज का \omterm{def:g4:measure:area}{क्षेत्रफल}
$\frac{6 \times 4}{2} = 12$ cm$^2$ है --- यह चित्र
\cref{ch:g6:measure} के लिए याद रखने लायक़ है, जहाँ यह सूत्र बन जाता है।
\end{example}

\begin{omfigure}
\begin{tikzpicture}[scale=0.62]
  \draw[gray!50, thin] (0,0) grid (6,4);
  \draw[very thick] (0,0) rectangle (6,4);
  \fill[omDef!25] (0,0) -- (6,0) -- (0,4) -- cycle;
  \draw[very thick, omDef] (6,0) -- (0,4);
  \node[font=\small] at (1.7,1.2) {आधा};
  \node[font=\small] at (4.3,2.8) {आधा};
\end{tikzpicture}

{\small विकर्ण $6 \times 4$ के आयत को $12$-$12$ खानों वाले दो
बराबर त्रिभुजों में काट देता है।}
\end{omfigure}

\section{अभ्यास}

\begin{exercise}[$\star$]\label{exo:g5:areas:1}
$8$ cm $\times$ $5$ cm के आयत का \omterm{def:g4:measure:area}{क्षेत्रफल} और \omterm{def:g4:measure:perimeter}{परिमाप} निकालो।
कौन-सा उत्तर cm में है, कौन-सा cm$^2$ में?
\end{exercise}

\begin{exercise}[$\star$]\label{exo:g5:areas:2}
\omterm{def:g4:measure:area}{क्षेत्रफल} निकालो: $9$ cm भुजा वाला वर्ग; $12$ m $\times$ $7$ m
का आयत; $4.5$ cm $\times$ $6$ cm का आयत।
\end{exercise}

\begin{exercise}[$\star$]\label{exo:g5:areas:3}
एक आयत का \omterm{def:g4:measure:area}{क्षेत्रफल} $63$ cm$^2$ और लंबाई $9$ cm है। उसकी चौड़ाई
निकालो, फिर उसका \omterm{def:g4:measure:perimeter}{परिमाप}।
\end{exercise}

\begin{exercise}[$\star$]\label{exo:g5:areas:4}
जाल वाले काग़ज़ पर $24$ खानों के \omterm{def:g4:measure:area}{क्षेत्रफल} वाले दो अलग-अलग आयत
बनाओ, और दोनों के \omterm{def:g4:measure:perimeter}{परिमाप} निकालो। \omterm{def:g4:measure:area}{क्षेत्रफल} वही --- \omterm{def:g4:measure:perimeter}{परिमाप} भी वही?
\end{exercise}

\begin{exercise}[$\star$]\label{exo:g5:areas:5}
बदलो: $3$ m$^2$ को cm$^2$ में; \quad $50\,000$ cm$^2$ को m$^2$ में।
(\cref{ex:g5:areas:units} याद रखो: सौ-सौ के हिसाब से!)
\end{exercise}

\begin{exercise}[$\star$]\label{exo:g5:areas:6}
T-आकार की एक आकृति $2 \times 5$ की खड़ी पट्टी के ऊपर रखी
$8 \times 2$ की आड़ी पट्टी से बनी है (cm में)। दो आयत जोड़कर
उसका \omterm{def:g4:measure:area}{क्षेत्रफल} निकालो।
\end{exercise}

\begin{exercise}[$\star$]\label{exo:g5:areas:7}
एक तस्वीर का चौखटा: $30$ cm $\times$ $20$ cm का आयत, जिसमें से
$24$ cm $\times$ $14$ cm का आयताकार झरोखा काट दिया गया है।
चौखटे में कितने \omterm{def:g4:measure:area}{क्षेत्रफल} की लकड़ी लगी है?
\end{exercise}

\begin{exercise}[$\star$]\label{exo:g5:areas:8}
$8$ cm और $6$ cm भुजाओं वाले \omterm{def:g3:shapes:rightangle}{समकोण} त्रिभुज का \omterm{def:g4:measure:area}{क्षेत्रफल} निकालो
(आयत का \omterm{def:g3:division:half}{आधा}, \cref{ex:g5:areas:halfrect})।
\end{exercise}

\begin{exercise}[$\star$]\label{exo:g5:areas:9}
एक आयताकार क्यारी $7$ m गुणा $4$ m की है। बीज $2$ प्रति वर्ग
मीटर का है, और बाड़ $3$ प्रति मीटर। बीज का ख़र्च निकालो, फिर बाड़
का ख़र्च।
\end{exercise}

\begin{exercise}[$\star\star$]\label{exo:g5:areas:10}
एक गलियारे का फ़र्श $12$ m लंबा और $2$ m चौड़ा है, और उस पर
$50$ cm भुजा वाली वर्गाकार टाइलें लगानी हैं। कितनी टाइलें
चाहिए? (पहले मात्रक बदलो, या हर दिशा में टाइलें गिनो।)
\end{exercise}

\begin{exercise}[$\star\star$]\label{exo:g5:areas:11}
$4$ cm $\times$ $3$ cm के आयत की भुजाएँ दुगुनी कर दो। उसके
\omterm{def:g4:measure:perimeter}{परिमाप} का क्या होता है? \omterm{def:g4:measure:area}{क्षेत्रफल} का? (पहले और बाद, दोनों
निकालो --- दोनों उत्तर अलग हैं!)
\end{exercise}
